\section{摩押之預言}
\subsection{摩押背景}

\begin{figure}
    \centering
    \begin{minipage}{0.35\textwidth}
        \centering
        \includegraphics[width=0.9\textwidth]{ZongJie/figs/MoYa}
    \end{minipage}\hfill
    \begin{minipage}{0.65\textwidth}
        \centering
        \includegraphics[width=1\textwidth]{ZongJie/figs/MOAB} 
    \end{minipage}
\end{figure}


\subsubsection{羅得大女兒生摩押 - 创19:30-38}
羅得因為怕住在瑣珥，就同他兩個女兒從瑣珥上去，住在山裡。他和兩個女兒住在一個洞裡。 31 大女兒對小女兒說：「我們的父親老了，地上又無人按著世上的常規進到我們這裡。 32 來，我們可以叫父親喝酒，與他同寢，這樣我們好從他存留後裔。」 33 於是那夜她們叫父親喝酒，大女兒就進去和她父親同寢。她幾時躺下，幾時起來，父親都不知道。 34 第二天，大女兒對小女兒說：「我昨夜與父親同寢，今夜我們再叫他喝酒，你可以進去與他同寢，這樣我們好從父親存留後裔。」 35 於是，那夜她們又叫父親喝酒，小女兒起來與她父親同寢。她幾時躺下，幾時起來，父親都不知道。 36 這樣，羅得的兩個女兒都從她父親懷了孕。 37 大女兒生了兒子，給他起名叫摩押，就是現今摩押人的始祖。 38 小女兒也生了兒子，給他起名叫便亞米，就是現今亞捫人的始祖。

\subsubsection{亞摩利人将摩押人赶出 - 民21:26-30}
摩押的发源地瑣珥——位于死海的东南边界——开始，他们逐渐扩展到约旦东部的地区。出埃及前不久，好战的亚摩利人在他们的王西宏的带领下越过约旦河，将摩押人赶出亚嫩河谷和雅博河之间的地区，使希实本成为他们的首都。


以色列人差遣使者去見亞摩利人的王西宏，說： 22 「求你容我們從你的地經過。我們不偏入田間和葡萄園，也不喝井裡的水，只走大道，直到過了你的境界。」 23 西宏不容以色列人從他的境界經過，就招聚他的眾民出到曠野，要攻擊以色列人，到了\textcolor{red}{雅雜}與以色列人爭戰。 24 以色列人用刀殺了他，得了他的地，從亞嫩河到雅博河，直到亞捫人的境界，因為亞捫人的境界多有堅壘。 

25 以色列人奪取這一切的城邑，也住亞摩利人的城邑，就是\textcolor{red}{希實本}與\textcolor{red}{希實本}的一切鄉村。 26 這\textcolor{red}{希實本}是亞摩利王西宏的京城，西宏曾與摩押的先王爭戰，從他手中奪取了全地，直到亞嫩河。 27 所以那些作詩歌的說：「你們來到\textcolor{red}{希實本}，願西宏的城被修造，被建立！ 28 因為有火從\textcolor{red}{希實本}發出，有火焰出於西宏的城，燒盡摩押的亞珥和亞嫩河丘壇的祭司。 29 摩押啊，你有禍了！\textcolor{red}{基抹}的民哪，你們滅亡了！\textcolor{red}{基抹}的男子逃奔，女子被擄，交付亞摩利的王西宏。 30 我們射了他們，\textcolor{red}{希實本}直到\textcolor{red}{底本}盡皆毀滅。我們使地變成荒場，直到\textcolor{red}{挪法}。這\textcolor{red}{挪法}直延到\textcolor{red}{米底巴}。」 31 這樣，以色列人就住在亞摩利人之地。 32 摩西打發人去窺探\textcolor{red}{雅謝}，以色列人就占了\textcolor{red}{雅謝}的鎮市，趕出那裡的亞摩利人。




\subsubsection{摩押王巴勒召巴蘭來咒詛以色列人 - 民22}
在出埃及期间，以色列人没有经过摩押，而是经过"旷野"向东，最终到达亚嫩河以北的地方。【2 以色列人向亞摩利人所行的一切事，西撥的兒子巴勒都看見了。 3 摩押人因以色列民甚多，就大大懼怕，心內憂急， 4 對米甸的長老說：「現在這眾人要把我們四圍所有的一概舔盡，就如牛舔盡田間的草一般。」那時西撥的兒子巴勒做\textcolor{red}{摩押王}。 5 他差遣使者往大河邊的\textcolor{red}{毗奪}去，到比珥的兒子巴蘭本鄉那裡，召巴蘭來，說：「有一宗民從埃及出來，遮滿地面，與我對居。 6 這民比我強盛，現在求你來為我咒詛他們，或者我能得勝，攻打他們，趕出此地。因為我知道，你為誰祝福，誰就得福；你咒詛誰，誰就受咒詛。」7 \textcolor{red}{摩押}的長老和米甸的長老手裡拿著卦金，到了巴蘭那裡，將巴勒的話都告訴了他。】(民22:2-7)。

【巴勒聽見巴蘭來了，就往\textcolor{red}{摩押}京城去迎接他。這城是在邊界上，在亞嫩河旁。 37 巴勒對巴蘭說：「我不是急急地打發人到你那裡去召你嗎？你為何不到我這裡來呢？我豈不能使你得尊榮嗎？」 38 巴蘭說：「我已經到你這裡來了。現在我豈能擅自說什麼呢？神將什麼話傳給我，我就說什麼。」 39 巴蘭和巴勒同行，來到\textcolor{red}{基列胡瑣}。 40 巴勒宰了牛羊，送給巴蘭和陪伴的使臣。41 到了早晨，巴勒領巴蘭到巴力的高處，巴蘭從那裡觀看以色列營的邊界。】 (民22:36-41)

\subsubsection{进入迦南地之前在\textcolor{red}{摩押}平原建营 - 民22}
以色列人起行，在摩押平原，約旦河東，對著耶利哥安營(民22:1)。

以上是摩西在約旦河東對著耶利哥的摩押平原所分給他們的產業。 33 只是利未支派，摩西沒有把產業分給他們，耶和華以色列的神是他們的產業，正如耶和華所應許他們的。(约13:31-32)

\subsubsection{摩西葬在\textcolor{red}{摩押}平原 - (申34:1-6)}
摩西從\textcolor{red}{摩押}平原登尼波山，上了那與耶利哥相對的毗斯迦山頂。耶和華把基列全地直到但， 2 拿弗他利全地，以法蓮、瑪拿西的地，猶大全地直到西海， 3 南地和棕樹城耶利哥的平原，直到瑣珥，都指給他看。 4 耶和華對他說：「這就是我向亞伯拉罕、以撒、雅各起誓應許之地，說：『我必將這地賜給你的後裔。』現在我使你眼睛看見了，你卻不得過到那裡去。」 5 於是，耶和華的僕人摩西死在摩押地，正如耶和華所說的。 6 耶和華將他埋葬在\textcolor{red}{摩押地}，伯毗珥對面的谷中，只是到今日沒有人知道他的墳墓。 7 摩西死的時候年一百二十歲，眼目沒有昏花，精神沒有衰敗。 8 以色列人在\textcolor{red}{摩押平原}為摩西哀哭了三十日，為摩西居喪哀哭的日子就滿了。



\subsubsection{摩押王伊矶伦为以笏所杀 - 士師記 3:15-30}
以色列人呼求耶和華的時候，耶和華就為他們興起一位拯救者，就是便雅憫人基拉的兒子以笏。他是左手便利的。以色列人託他送禮物給摩押王伊磯倫。 16 以笏打了一把兩刃的劍，長一肘，帶在右腿上衣服裡面。 17 他將禮物獻給摩押王伊磯倫，原來伊磯倫極其肥胖。 18 以笏獻完禮物，便將抬禮物的人打發走了， 19 自己卻從靠近吉甲鑿石之地回來，說：「王啊，我有一件機密事奏告你。」王說：「迴避吧！」於是左右侍立的人都退去了。 20 以笏來到王面前，王獨自一人坐在涼樓上。以笏說：「我奉神的命報告你一件事。」王就從座位上站起來。 21 以笏便伸左手，從右腿上拔出劍來，刺入王的肚腹， 22 連劍把都刺進去了。劍被肥肉夾住，他沒有從王的肚腹拔出來，且穿通了後身。 23 以笏就出到遊廊，將樓門盡都關鎖。24 以笏出來之後，王的僕人到了，看見樓門關鎖，就說：「他必是在樓上大解。」 25 他們等煩了，見仍不開樓門，就拿鑰匙開了，不料，他們的主人已死，倒在地上。26 他們耽延的時候，以笏就逃跑了，經過鑿石之地，逃到西伊拉。 27 到了，就在以法蓮山地吹角。以色列人隨著他下了山地，他在前頭引路， 28 對他們說：「你們隨我來，因為耶和華已經把你們的仇敵摩押人交在你們手中。」於是他們跟著他下去，把守約旦河的渡口，不容摩押一人過去。 29 那時擊殺了摩押人約有一萬，都是強壯的勇士，沒有一人逃脫。 30 這樣，摩押就被以色列人制伏了。國中太平八十年。


\subsubsection{扫罗王曾经打败摩押等诸族(撒上14:47)}
掃羅執掌以色列的國權，常常攻擊他四圍的一切仇敵，就是摩押人、亞捫人、以東人和瑣巴諸王，並非利士人。他無論往何處去，都打敗仇敵。


\subsubsection{大卫王和摩押等诸族(撒上14:47)}
大卫的祖先路得是摩押女子，大卫躲避扫罗的时候，曾求摩押王容他父亲住在摩押地(撒上22:3，4)；以后又攻打摩押人(犹太传统认为，大卫之所以攻打摩押人，是因为摩押人后来杀了大卫的父母。)，使他们或被杀，或归服而进贡(撒下8:2，诗60:8)。


\subsubsection{所罗门王和摩押 (撒上14:47)}
所罗门王的妃嫔中，有摩押的女子，所罗门王为她们的神建筑邱坛。所罗门王死后，


摩押复兴，直到以色列王暗利在位时，始再失败；到了亚哈王死后，摩押又复背叛(王下3:5)，结果为以色列王、犹大王和以东王所败，但不久又自立，时常侵犯以色列境；

\subsubsection{众先知皆曾预言摩押必受刑罚}
众先知皆曾预言摩押必受刑罚，必要覆亡。主前七百馀年，摩押进贡亚述。在约雅敬王年间，摩押军和亚扪军，以及外邦各军，齐来攻击约雅敬，毁灭犹大(王下24:2)。到了犹大人从巴比伦回归的时候，阿拉伯的拿巴底安族已占据了摩押地，摩押人的权势就日渐衰微。在圣经之外现存最后关于摩押的最后史料(cylinder B of Assurbanipal)，是摩押王Kaasḥalta臣服于亚述王亚述巴尼拔(主前669-631年在位)。从波斯帝国开始，摩押不再出现在史料中，只有他们的故土还被称为「摩押地」，但却被一波又一波来自北阿拉伯的部落占据。

   

\subsubsection{希實本 - 迦得,支派,利未支派米拉利子孫}
位於約但河東岸、鹽海東北部的高地上，西距耶路撒冷80公里。本為摩押人居住地，後被亞摩利王西宏以武力攻佔，立為首都。以色列人進軍迦南途中曾向西宏借道，遭拒絕後打敗西宏的軍隊，殺死所有男女老幼，奪取希實本全境。分配地業時，該城劃歸流便支派，因在流便屬地與迦得屬地的交界處，實際歸於迦得支派。而後又遵從耶和華旨意，將希實本、瑪哈念等四城和屬城的郊野給了利未支派米拉利的子孫為業。大約在西元前853年左右，摩押王米沙又從以色列人手中奪回希實本。而後亞捫佔據了該城。


\begin{table}[H]
\centering 
\begin{tabular}{p{2.cm}|p{14cm}|p{1.5cm}}\hline\hline
&以色列人奪取這一切的城邑，也住亞摩利人的城邑，就是\textcolor{red}{希實本}與\textcolor{red}{希實本}的一切鄉村。 26 這\textcolor{red}{希實本}是亞摩利王西宏的京城，西宏曾與摩押的先王爭戰，從他手中奪取了全地，直到亞嫩河。 27 所以那些作詩歌的說：「你們來到\textcolor{red}{希實本}，願西宏的城被修造，被建立！ 28 因為有火從\textcolor{red}{希實本}發出，有火焰出於西宏的城，燒盡摩押的亞珥和亞嫩河丘壇的祭司。 29 摩押啊，你有禍了！基抹的民哪，你們滅亡了！基抹的男子逃奔，女子被擄，交付亞摩利的王西宏。 30 我們射了他們，\textcolor{red}{希實本}直到底本盡皆毀滅。我們使地變成荒場，直到挪法。這挪法直延到米底巴。」...以色列人轉回，向巴珊去。巴珊王噩和他的眾民都出來，在以得來與他們交戰。 34 耶和華對摩西說：「不要怕他。因我已將他和他的眾民，並他的地，都交在你手中。你要待他像從前待住\textcolor{red}{希實本}的亞摩利王西宏一般。」&民21:25-34\\\hline
&魯本子孫和迦得子孫的牲畜極其眾多。他們看見雅謝地和基列地是可牧放牲畜之地， 2 就來見摩西和祭司以利亞撒並會眾的首領，說： 3 「亞他錄、底本、雅謝、寧拉、\textcolor{red}{希實本}、以利亞利、示班、尼波、比穩， 4 就是耶和華在以色列會眾前面所攻取之地，是可牧放牲畜之地，你僕人也有牲畜。」 5 又說：「我們若在你眼前蒙恩，求你把這地給我們為業，不要領我們過約旦河。」&民32:1-5\\\hline
&摩西將亞摩利王西宏的國和巴珊王噩的國，連那地和周圍的城邑，都給了迦得子孫和魯本子孫，並約瑟的兒子瑪拿西半個支派。...魯本子孫建造\textcolor{red}{希實本}、以利亞利、基列亭、 38 尼波、巴力免、西比瑪（尼波、巴力免，名字是改了的），又給他們所建造的城另起別名。&民32:37\\\hline
&出埃及第四十年十一月初一日，摩西照耶和華藉著他所吩咐以色列人的話，都曉諭他們。 4 那時，他已經擊殺了住希實本的亞摩利王西宏和住以得來、亞斯他錄的巴珊王噩。&申1:4\\\hline
&你們起來前往，過亞嫩谷。我已將亞摩利人\textcolor{red}{希實本}王西宏和他的地交在你手中，你要與他爭戰，得他的地為業。 25 從今日起，我要使天下萬民聽見你的名聲都驚恐懼怕，且因你發顫傷慟。我從基底莫的曠野差遣使者去見\textcolor{red}{希實本}王西宏，用和睦的話說： 27 『求你容我從你的地經過，只走大道，不偏左右。 28 你可以賣糧給我吃，也可以賣水給我喝，只要容我步行過去， 29 就如住西珥的以掃子孫和住亞珥的摩押人待我一樣，等我過了約旦河，好進入耶和華我們神所賜給我們的地。』 30 但\textcolor{red}{希實本}王西宏不容我們從他那裡經過，因為耶和華你的神使他心中剛硬，性情頑梗，為要將他交在你手中，像今日一樣。 31 耶和華對我說：『從此起首，我要將西宏和他的地交給你。你要得他的地為業！』 32 那時，西宏和他的眾民出來攻擊我們，在\textcolor{red}{雅雜}與我們交戰。 33 耶和華我們的神將他交給我們，我們就把他和他的兒子並他的眾民，都擊殺了。&申2:24-33\\\hline
戰敗巴珊王噩&我們轉回，向巴珊去。巴珊王噩和他的眾民都出來，在以得來與我們交戰。 2 耶和華對我說：『不要怕他，因我已將他和他的眾民並他的地，都交在你手中。你要待他像從前待住\textcolor{red}{希實本}的亞摩利王西宏一樣。』....我們將這些都毀滅了，像從前待\textcolor{red}{希實本}王西宏一樣，把有人煙的各城，連女人帶孩子，盡都毀滅。&申3:2，6\\\hline
摩西在以色列人前陳明律法&在約旦河東伯毗珥對面的谷中，在住\textcolor{red}{希實本}亞摩利王西宏之地。這西宏是摩西和以色列人出埃及後所擊殺的， 47 他們得了他的地，又得了巴珊王噩的地，就是兩個亞摩利王在約旦河東向日出之地，&申4:46\\\hline
耶和華在摩押地立約的話&你們來到這地方，\textcolor{red}{希實本}王西宏、巴珊王噩都出來與我們交戰，我們就擊殺了他們， 8 取了他們的地給魯本支派、迦得支派和瑪拿西半支派為業......&申29:7-8\\\hline
\end{tabular}
\caption{希實本}
\end{table}


\begin{table}[H]
\centering 
\begin{tabular}{p{2.cm}|p{14cm}|p{1.5cm}}\hline\hline
&將列國之地照份賜給他們，他們就得了西宏之地、\textcolor{red}{希實本}王之地和巴珊王噩之地。&尼9:22\\\hline
&基遍人回答說：「僕人從極遠之地而來，是因聽見耶和華你神的名聲和他在埃及所行的一切事， 10 並他向約旦河東的兩個亞摩利王，就是\textcolor{red}{希實本}王西宏和在亞斯他錄的巴珊王噩一切所行的事。&書9:9-10\\\hline
&以色列人在約旦河外向日出之地擊殺二王，得他們的地，就是從亞嫩谷直到黑門山，並東邊的全亞拉巴之地。 2 這二王，有住\textcolor{red}{希實本}亞摩利人的王西宏。&書12:1-2\\\hline
&瑪拿西那半支派和魯本、迦得二支派已經受了產業...在\textcolor{red}{希實本}做王亞摩利王西宏的諸城，直到亞捫人的境界；...\textcolor{red}{希實本}並屬\textcolor{red}{希實本}平原的各城，底本，巴末巴力，伯巴力勉， 18 雅雜，基底莫，米法押，摩西按著迦得支派的宗族分給他們產業。 25 他們的境界是\textcolor{red}{雅謝}和基列的各城，並亞捫人的一半地，直到拉巴前的亞羅珥； 26 從\textcolor{red}{希實本}到拉抹米斯巴和比多寧，又從瑪哈念到底璧的境界， 27 並谷中的伯亞蘭、伯寧拉、疏割、撒分，就是\textcolor{red}{希實本}王西宏國中的餘地；...以上是迦得人按著宗族所得為業的諸城，並屬城的村莊。&書13:10,17,26-27\\\hline
&其餘利未支派米拉利子孫，從西布倫支派的地業中所得的，就是約念和屬城的郊野，加珥他和屬城的郊野， 35 丁拿和屬城的郊野，拿哈拉和屬城的郊野，共四座城。 36 又從魯本支派的地業中給了他們比悉和屬城的郊野，\textcolor{red}{雅雜}和屬城的郊野， 37 基底莫和屬城的郊野，米法押和屬城的郊野，共四座城。 38 又從迦得支派的地業中，將基列的拉末，就是誤殺人的逃城和屬城的郊野，給了他們，又給他們瑪哈念和屬城的郊野， 39 \textcolor{red}{希實本}和屬城的郊野，\textcolor{red}{雅謝}和屬城的郊野，共四座城。 40 其餘利未支派的人，就是米拉利的子孫，按著宗族拈鬮所得的，共十二座城。&書21:38-39\\\hline
&以色列人打發使者去見亞摩利王西宏，就是 \textcolor{red}{希實本}的王，對他說：「求你容我們從你的地經過，往我們自己的地方去。」 20 西宏卻不信服以色列人，不容他們經過他的境界，乃招聚他的眾民在\textcolor{red}{雅雜}安營，與以色列人爭戰。 21 耶和華以色列的神將西宏和他的眾民都交在以色列人手中，以色列人就擊殺他們，得了亞摩利人的全地， 22 從亞嫩河到雅博河，從曠野直到約旦河。 23 耶和華以色列的神在他百姓以色列面前趕出亞摩利人，你竟要得他們的地嗎？ 24 你的神基抹所賜你的地，你不是得為業嗎？耶和華我們的神在我們面前所趕出的人，我們就得他的地。 25 難道你比摩押王西撥的兒子巴勒還強嗎？他曾與以色列人爭競，或是與他們爭戰嗎？ 26 以色列人住 \textcolor{red}{希實本}和屬希實本的鄉村，亞羅珥和屬亞羅珥的鄉村，並沿亞嫩河的一切城邑，已經有三百年了。在這三百年之內，你們為什麼沒有取回這些地方呢？&士11:19-26\\\hline
&又在迦得支派的地中得了基列的拉末與其郊野，瑪哈念與其郊野， 81 \textcolor{red}{希實本}與其郊野，\textcolor{red}{雅謝}與其郊野。&代上6:80-81\\\hline
&一夜之間，摩押的亞珥變為荒廢，歸於無有。...\textcolor{red}{希實本}和以利亞利悲哀的聲音達到\textcolor{red}{雅雜}，所以摩押帶兵器的高聲喊嚷，人心戰兢。&賽15:4\\\hline
&因為\textcolor{red}{希實本}的田地和西比瑪的葡萄樹都衰殘了。列國的君主折斷其上美好的枝子，這枝子長到\textcolor{red}{雅謝}，延到曠野；嫩枝向外探出，直探過鹽海。 9 「因此，我要為西比瑪的葡萄樹哀哭，與雅謝人哀哭一樣。\textcolor{red}{希實本}、以利亞利啊，我要以眼淚澆灌你，因為有交戰呐喊的聲音臨到你夏天的果子，並你收割的莊稼。&賽16:8-9\\\hline
&摩押不再被稱讚，有人在\textcolor{red}{希實本}設計謀害她說：『來吧！我們將她剪除，不再成國。』瑪得緬哪，你也必默默無聲，刀劍必追趕你。....肥田和摩押地的歡喜快樂都被奪去，我使酒榨的酒絕流，無人踹酒歡呼，那歡呼卻變為仇敵的呐喊。 34 \textcolor{red}{希實本}人發的哀聲達到以利亞利，直達到雅雜，從瑣珥達到何羅念，直到伊基拉施利施亞，因為寧林的水必然乾涸。....躲避的人無力站在\textcolor{red}{希實本}的影下，因為有火從\textcolor{red}{希實本}發出，有火焰出於西宏的城，燒盡摩押的角和鬨嚷人的頭頂。&耶48:2,33-34,45\\\hline
&論亞捫說，「\textcolor{red}{希實本}哪，你要哀號，因為愛地變為荒場」&耶49:3\\\hline
&你的眼目像\textcolor{red}{希實本}巴特拉併門旁的水池。&歌7:4\\\hline
\end{tabular}
\caption{希實本}
\end{table}

\subsubsection{底本 （底們） - Dimon}
 亞羅珥-1   (Aroer-1) [詳解編號 : 0565-1]
廢墟 Ruins，貧乏的 Destitute，包圍 Enclosed
●Arair，位於底本東南約 5 公里，希實本之南約 36 公里，亞嫩河的北岸，距死海約 24 公里，曾是摩押一個設了防的重鎮，並有東方的別是巴之稱。

民32:34 是迦得的子孫所建造的八個堅固城之一。

申2:36，3:12，4:48，書13:9，16 以色列人奪取了亞摩利人的王西宏的國，將亞嫩谷邊的亞羅珥，和谷中的城等給了流便支派為地業。

書12:2 亞摩利人的王西宏，他所管之地，是從亞嫩谷邊的亞羅珥和谷中的城，並基列一半。

士11:26 士師耶弗他對亞捫人說，以色列人住希實本，亞羅珥，並沿亞嫩河的一切城邑，已經有三百年了。

王下10:33 在以色列王耶戶年間，耶和華才割裂以色列國，使哈薛攻擊以色列的境界，乃是約但河東，基列全地，從靠近亞嫩谷邊的亞羅珥起，就是基列和巴珊的迦得人，流便人，瑪拿西人之地。

代上5:8 流便的後代，約珥所住的地方是從亞羅珥直到尼波和巴力免。

代上11:44 大衛的勇士中，有亞羅珥人沙瑪和耶利等兩個人。

\subsubsection{拉末 - Ramoth}
拉末-1   (Ramoth-1) [詳解編號 : 0634-1]
高地 High places，Heights，珊瑚 Coral
●T.Ramith，位於伯珊以東約 51 公里，以得來西南約 11 公里，是從東方進入基列的孔道，君王大道也通過此城，自古就是一個重要的軍事據點，現今有兩條路幹道從其北經過。在聖經中多稱作基列的拉末 Ramoth-Gilead。

申4:43，書20:8，21:38 摩西在約但河之東分定三座逃城，其中為迦得人分定基列的拉末為逃城。

王上4:13 所羅門王的十二個行政區中的第六區，由便基別管理，其首府是基列的拉末。

王上22:3，代下18:2 以色列王亞哈聯合猶大王約沙法，一同攻取被亞蘭人所佔領的基列拉末，結果大敗，亞哈王被箭射死。

王下8:25，代下22:5 以色列王約蘭聯合猶大王亞哈謝，前往基列的拉末去，與亞蘭王哈薛爭戰，結果大敗，約蘭受傷後死亡。

王下9:1 先知以利沙吩咐他的一個門徒到基列的拉末去，膏耶戶作以色列王。
代上6:80 利未支派的米拉利族從迦得支派的地業中得了基列的拉末，瑪哈念，希實本和雅謝等城及其郊野為業。



\subsubsection{吉珥哈列設 (Kir-haraseth)}
摩押的一個重要防城，有說是摩押的首都。在鹽海南端的東面、由大馬色通往以旬迦別的南北大道上。以色列第九位王約蘭聯合猶大王約沙法及以東王共同進攻摩押，大獲全勝，殺人摩押境內，破壞摩押的一切城邑、水泉、樹木等，僅剩下吉珥哈列設的石牆。先知以賽亞和耶利米都預言摩押的災難將臨，人們要為吉珥哈列設歎息憂傷。有說 「吉珥」是其簡稱。參見該條2。(王下3:24-25；賽16:6-12；耶48 :31)該城即今距死海南端約18公里的克拉克。


吉珥哈列設   (Kir-hareseth) [詳解編號 : 0250]
陶器之城 City of pottery, 新城 New city
●el-Kerak，是約旦今日的一個大城，位於底本之南約 36 公里，亞珥之南約 11 公里，漓三半島之西約 20 公里，城建在一高地上，標高約 1100 公尺，四週又有山嶺圍繞，故難攻易守，曾經是摩押王國的京城，其城堡和城牆都築造得十分堅固雄壯，自今仍保留了大部份。昔時的君王大道經過此處，今日則有一條東西向的公路通過，所以自古就是一個極具商業和軍事價值的重鎮。

王下3:25 以色列王約蘭，猶大王約沙法和以東王聯合攻打摩押，摩押王在吉珥哈列設的石牆上，將那應接續他作王的長子，在城上獻為燔祭。

賽16:7，耶48:31 先知預言摩押因狂傲，吉珥哈列設必將遭災。

※同賽 15:1 中的 ( 摩押的 ) 基珥 Kir 0912，是其簡稱，曾被先知所咒詛。

\subsubsection{底本 （底們） - Dimon}
摩押地一城市。先知以賽亞指斥底們是摩押諸惡城之首，必受耶和華的懲罰。(賽15:9)。有說即摩押地的「底本」。「底們」與「底本」諧音，字根的意思是「血」，預言底本血腥的命運。參見「底本1」。其遺址在今拉伊拜市西北約5公里處的底穆奈廢墟。

位於鹽海以東、亞嫩河北面約5公里處。原屬摩押人，亞摩利王西宏戰勝摩押人後據為己有。以色列人進軍迦南途中在此安營，向西宏借道遭拒絕後，打敗西宏，佔領該地。分配地業時，劃歸迦得的子孫(民33-34 說，迦得支派建造了底本、亞他錄、雅謝等城市；書13:17則說劃給流便的子孫為業) ，所以底本又稱「底本迦得」。以色列南北分裂時代，國勢衰落，底本又為摩押控制。先知以賽亞等指斥底們是摩押諸惡城之首。「底們」與「底本」諧音，字根的意思是「血」，預言底本血腥的命運。(民21:30；32:3，33-34，45-46；賽15:2，9；耶48:18，22)今名「底班」，在約旦境內。

\subsubsection{雅雜（Jahaz, Jahazah) - 流便支派,利未支派米拉利子孫}
鹽海東面一城鎮。原為亞摩利人居住地。以色列人進軍迦南途中，在此打敗亞摩利人希實本王西宏，奪取該地和希實本的所有城市、牲畜與財產。摩西分配地業時，劃歸流便的子孫。後遵從耶和華旨意，將雅雜等四城和屬城的郊野給了利未支派米拉利的子孫為業。許多年後，摩押人戰勝以色列王米沙，奪得雅雜。確址不詳。

\begin{table}[H]
\centering 
\begin{tabular}{p{2.cm}|p{14cm}|p{1.cm}}\hline\hline
&以色列人差遣使者去見亞摩利人的王西宏，說： 22 「求你容我們從你的地經過。我們不偏入田間和葡萄園，也不喝井裡的水，只走大道，直到過了你的境界。」 23 西宏不容以色列人從他的境界經過，就招聚他的眾民出到曠野，要攻擊以色列人，到了\textcolor{red}{雅雜}與以色列人爭戰。 24 以色列人用刀殺了他，得了他的地，從亞嫩河到雅博河，直到亞捫人的境界，因為亞捫人的境界多有堅壘。 25 以色列人奪取這一切的城邑，也住亞摩利人的城邑，就是希實本與希實本的一切鄉村。 26 這希實本是亞摩利王西宏的京城，西宏曾與摩押的先王爭戰，從他手中奪取了全地，直到亞嫩河。 27 所以那些作詩歌的說：「你們來到希實本，願西宏的城被修造，被建立！ 28 因為有火從希實本發出，有火焰出於西宏的城，燒盡摩押的亞珥和亞嫩河丘壇的祭司。 29 摩押啊，你有禍了！基抹的民哪，你們滅亡了！基抹的男子逃奔，女子被擄，交付亞摩利的王西宏。 30 我們射了他們，希實本直到底本盡皆毀滅。我們使地變成荒場，直到挪法。這挪法直延到米底巴。」 31 這樣，以色列人就住在亞摩利人之地。 32摩西打發人去窺探 \textcolor{red}{雅謝}，以色列人就占了雅謝的鎮市，趕出那裡的亞摩利人。以色列人轉回，向巴珊去。巴珊王噩和他的眾民都出來，在以得來與他們交戰。 34 耶和華對摩西說：「不要怕他。因我已將他和他的眾民，並他的地，都交在你手中。你要待他像從前待住希實本的亞摩利王西宏一般。」於是他們殺了他和他的眾子,並他的眾民,沒有留下一個,就得了他的地。&民21:32\\\hline


&魯本子孫和迦得子孫的牲畜極其眾多。他們看見雅謝地和基列地是可牧放牲畜之地， 2 就來見摩西和祭司以利亞撒並會眾的首領說： 3 「亞他錄,底本,\textcolor{red}{雅謝},寧拉,\textcolor{red}{希實本},以利亞利,示班,尼波,比穩, 4 就是耶和華在以色列會眾前面所攻取之地是可牧放牲畜之地,你僕人也有牲畜。」 5 又說：「我們若在你眼前蒙恩，求你把這地給我們為業,不要領我們過約旦河。」&民32:3，37\\\hline


&摩西將亞摩利王西宏的國和巴珊王噩的國，連那地和周圍的城邑，都給了迦得子孫和魯本子孫，並約瑟的兒子瑪拿西半個支派。 34 迦得子孫建造底本、亞他錄、亞羅珥、 35 亞他錄朔反、\textcolor{red}{雅謝}、約比哈、 36 伯寧拉、伯哈蘭，都是堅固城，他們又壘羊圈。&民32:34-35\\\hline

摩西在約旦河東曠野向以色列人說&那時，西宏和他的眾民出來攻擊我們，在\textcolor{red}{雅雜}與我們交戰。 33 耶和華我們的神將他交給我們，我們就把他和他的兒子並他的眾民，都擊殺了。 34 我們奪了他的一切城邑，將有人煙的各城，連女人帶孩子，盡都毀滅，沒有留下一個。 35 唯有牲畜和所奪的各城，並其中的財物，都取為自己的掠物。&申2:32-35\\\hline
&摩西按著迦得支派的宗族分給他們產業。 25 他們的境界是\textcolor{red}{雅謝}和基列的各城，並亞捫人的一半地，直到拉巴前的亞羅珥； ...就是希實本王西宏國中的餘地；以及約旦河與靠近約旦河的地，直到基尼烈海的極邊，都在約旦河東。&書13:25-26\\\hline
&其餘利未支派米拉利子孫，從西布倫支派的地業中所得的，就是約念和屬城的郊野，加珥他和屬城的郊野， 35 丁拿和屬城的郊野，拿哈拉和屬城的郊野，共四座城。 36 又從魯本支派的地業中給了他們比悉和屬城的郊野，\textcolor{red}{雅雜}和屬城的郊野， 37 基底莫和屬城的郊野，米法押和屬城的郊野，共四座城。 38 又從迦得支派的地業中，將基列的拉末，就是誤殺人的逃城和屬城的郊野，給了他們，又給他們瑪哈念和屬城的郊野， 39 希實本和屬城的郊野，\textcolor{red}{雅謝}和屬城的郊野，共四座城。 40 其餘利未支派的人，就是米拉利的子孫，按著宗族拈鬮所得的，共十二座城。&書21:39\\\hline
耶弗他打發使者去見亞捫人的王&以色列人打發使者去見亞摩利王西宏，就是希實本的王，對他說：「求你容我們從你的地經過，往我們自己的地方去。」 20 西宏卻不信服以色列人，不容他們經過他的境界，乃招聚他的眾民在\textcolor{red}{雅雜}安營，與以色列人爭戰。 21 耶和華以色列的神將西宏和他的眾民都交在以色列人手中，以色列人就擊殺他們，得了亞摩利人的全地，&士11:20\\\hline
論摩押的默示&一夜之間，摩押的亞珥變為荒廢，歸於無有。一夜之間，摩押的基珥變為荒廢，歸於無有。 2 他們上巴益，又往底本，到高處去哭泣。摩押人因尼波和米底巴哀號，各人頭上光禿，鬍鬚剃淨。 3 他們在街市上都腰束麻布，在房頂上和寬闊處俱各哀號，眼淚汪汪。 4 希實本和以利亞利悲哀的聲音達到\textcolor{red}{雅雜}，所以摩押帶兵器的高聲喊嚷，人心戰兢。&賽15:4\\\hline
耶和華說&摩押因毀壞蒙羞,你們要哀號呼喊,要在亞嫩旁報告說：『摩押變為荒場！』 21 刑罰臨到平原之地的何倫,\textcolor{red}{雅雜}...波斯拉和摩押地遠近所有的城邑。 摩押的角砍斷了,摩押的膀臂折斷了&耶48:21\\\hline
\end{tabular}
\caption{雅雜（Jahaz, Jahazah)}
\end{table}







\subsection{以賽亞書}
\subsubsection{預示摩押必覆亡 15:1-9}


論摩押的默示。一夜之間，摩押的亞珥變為荒廢，歸於無有。一夜之間，摩押的基珥變為荒廢，歸於無有。 2 他們上巴益，又往底本，到高處去哭泣。摩押人因尼波和米底巴哀號，各人頭上光禿，鬍鬚剃淨。 3 他們在街市上都腰束麻布，在房頂上和寬闊處俱各哀號，眼淚汪汪。 4 希實本和以利亞利悲哀的聲音達到雅雜，所以摩押帶兵器的高聲喊嚷，人心戰兢。 5 「我心為摩押悲哀，她的貴胄 (逃民)逃到瑣珥，到伊基拉、施利施亞。他們上魯希坡隨走隨哭，在何羅念的路上因毀滅舉起哀聲。 6 因為寧林的水成為乾涸，青草枯乾，嫩草滅沒，青綠之物一無所有。 7 因此，摩押人所得的財物和所積蓄的都要運過柳樹河。 8 哀聲遍聞摩押的四境，哀號的聲音達到以基蓮，哀號的聲音達到比珥以琳。 9 底們的水充滿了血，我還要加增底們的災難，叫獅子來追上摩押逃脫的民和那地上所餘剩的人。」

\subsubsection{勸摩押納貢於錫安之君 16:1-16}
你們當將羊羔奉給那地掌權的，從西拉往曠野，送到錫安城(女子)的山。 2 摩押的居民 (女子)在亞嫩渡口，必像遊飛的鳥，如拆窩的雛。 3 「求你獻謀略，行公平，使你的影子在午間如黑夜，隱藏被趕散的人，不可顯露逃民。 4 求你容我這被趕散的人和你同居；至於摩押，求你做她的隱密處，脫離滅命者的面。」勒索人的歸於無有，毀滅的事止息了，欺壓人的從國中除滅了。 5 必有寶座因慈愛堅立，必有一位誠誠實實坐在其上，在大衛帳幕中施行審判，尋求公平，速行公義。

%\subsubsection{摩押狂傲必遭災罰，16:6-7}
我們聽說摩押人驕傲，是極其驕傲，聽說她狂妄、驕傲、憤怒，她誇大的話是虛空的。 7 因此，摩押人必為摩押哀號，人人都要哀號。你們摩押人要為吉珥哈列設的葡萄餅哀嘆，極其憂傷。

%\subsubsection{先知哀嘆摩押人到他聖所祈禱也不蒙應允 16:8-12}
因為希實本的田地和西比瑪的葡萄樹都衰殘了。列國的君主折斷其上美好的枝子，這枝子長到雅謝，延到曠野；嫩枝向外探出，直探過鹽海。 9 「因此，我要為西比瑪的葡萄樹哀哭，與雅謝人哀哭一樣。希實本、以利亞利啊，我要以眼淚澆灌你，因為有交戰呐喊的聲音臨到你夏天的果子，並你收割的莊稼。從肥美的田中奪去了歡喜快樂，在葡萄園裡必無歌唱，也無歡呼的聲音；踹酒的在酒榨中不得踹出酒來，我使他歡呼的聲音止息。 11 因此我心腹為摩押哀鳴如琴，我心腸為吉珥哈列設也是如此。 12 摩押人朝見的時候，在高處疲乏，又到他聖所祈禱，也不蒙應允。」

%\subsubsection{預定摩押三年之內遭災 16:13-14}
這是耶和華從前論摩押的話。 14 但現在耶和華說：「三年之內，照雇工的年數，摩押的榮耀與她的群眾必被藐視，餘剩的人甚少無幾。」

%subsubsection{因賜救恩歡喜，摩押人之地必被踐踏 25:9-12}
到那日，人必說：「看哪，這是我們的神，我們素來等候他，他必拯救我們。這是耶和華，我們素來等候他，我們必因他的救恩歡喜快樂。」 10 耶和華的手必按在這山上。摩押人在所居之地必被踐踏，好像乾草被踐踏在糞池的水中。 11 她必在其中伸開手，好像洑水的伸開手洑水一樣，但耶和華必使她的驕傲和她手所行的詭計一併敗落。 12 耶和華使你城上的堅固高臺傾倒，拆平直到塵埃。

\subsubsection{摩押人在之地必被踐踏 25:10-12}
耶和華的手必按在這山上。摩押人在所居之地必被踐踏，好像乾草被踐踏在糞池的水中。 11 她必在其中伸開手，好像洑水的伸開手洑水一樣，但耶和華必使她的驕傲和她手所行的詭計一併敗落。 12 耶和華使你城上的堅固高臺傾倒，拆平直到塵埃。

\subsection{耶利米書}
\subsubsection{預言摩押必受懲罰 48:1-6}
論摩押。萬軍之耶和華以色列的神如此說：「尼波有禍了，因變為荒場。基列亭蒙羞被攻取，米斯迦蒙羞被毀壞。 2 摩押不再被稱讚，有人在希實本設計謀害她說：『來吧！我們將她剪除，不再成國。』瑪得緬哪，你也必默默無聲，刀劍必追趕你。 3 從何羅念有喊荒涼大毀滅的哀聲。 4 摩押毀滅了，她的孩童(家童)發哀聲，使人聽見。 5 人上魯希坡隨走隨哭，因為在何羅念的下坡聽見毀滅的哀聲。 6 你們要奔逃，自救性命，獨自居住，好像曠野的杜松。

\subsubsection{因自恃財貨 48:7-10}
「你因倚靠自己所做的和自己的財寶，必被攻取，基抹和屬他的祭司、首領也要一同被擄去。 8 行毀滅的必來到各城，並無一城得免。山谷必致敗落，平原必被毀壞，正如耶和華所說的。 9 要將翅膀給摩押，使她可以飛去。她的城邑必致荒涼，無人居住。 10 懶惰為耶和華行事的，必受咒詛；禁止刀劍不經血的，必受咒詛。

\subsubsection{因妄賴勇力 48:11-25}
「摩押自幼年以來常享安逸，如酒在渣滓上澄清，沒有從這器皿倒在那器皿裡，也未曾被擄去。因此她的原味尚存，香氣未變。」 12 耶和華說：「日子將到，我必打發倒酒的往她那裡去，將她倒出來，倒空她的器皿，打碎她的罈子。 13 摩押必因基抹羞愧，像以色列家從前倚靠伯特利的神羞愧一樣。 14 你們怎麼說『我們是勇士，是有勇力打仗的』呢？ 15 摩押變為荒場，敵人上去進了她的城邑，她所特選的少年人下去遭了殺戮。」這是君王，名為萬軍之耶和華說的。 16 「摩押的災殃臨近，她的苦難速速來到。 17 凡在她四圍的和認識她名的，你們都要為她悲傷，說：『那結實的杖和那美好的棍何竟折斷了呢！』 18 住在底本的民(女子)哪，要從你榮耀的位上下來，坐受乾渴，因毀滅摩押的上來攻擊你，毀壞了你的保障。 19 住亞羅珥的啊，要站在道旁觀望，問逃避的男人和逃脫的女人說：『是什麼事呢？』 20 摩押因毀壞蒙羞，你們要哀號呼喊，要在亞嫩旁報告說：『摩押變為荒場！』 21 刑罰臨到平原之地的何倫、雅雜、米法押、 22 底本、尼波、伯低比拉太音、 23 基列亭、伯迦末、伯米恩、 24 加略、波斯拉，和摩押地遠近所有的城邑。 25 摩押的角砍斷了，摩押的膀臂折斷了。」這是耶和華說的。

\subsubsection{因悖逆神輕侮主民 48:26-44}
「你們要使摩押沉醉，因她向耶和華誇大。她要在自己所吐之中打滾，又要被人嗤笑。 27 摩押啊，你不曾嗤笑以色列嗎？她豈是在賊中查出來的呢？你每逢提到她便搖頭。 28 摩押的居民哪，要離開城邑，住在山崖裡，像鴿子在深淵口上搭窩。 29 我們聽說摩押人驕傲，是極其驕傲；聽說她自高自傲，並且狂妄，居心自大。」 30 耶和華說：「我知道她的憤怒是虛空的，她誇大的話一無所成。 31 因此我要為摩押哀號，為摩押全地呼喊，人必為吉珥哈列設人嘆息。 32 西比瑪的葡萄樹啊，我為你哀哭，甚於雅謝人哀哭。你的枝子蔓延過海，直長到雅謝海，那行毀滅的已經臨到你夏天的果子和你所摘的葡萄。 33 肥田和摩押地的歡喜快樂都被奪去，我使酒榨的酒絕流，無人踹酒歡呼，那歡呼卻變為仇敵的呐喊(那歡呼卻不是歡呼)。 34 希實本人發的哀聲達到以利亞利，直達到雅雜，從瑣珥達到何羅念，直到伊基拉施利施亞，因為寧林的水必然乾涸。」 35 耶和華說：「我必在摩押地使那在丘壇獻祭的和那向她的神燒香的，都斷絕了。

「我心腹為摩押哀鳴如簫，我心腸為吉珥哈列設人也是如此，因摩押人所得的財物都滅沒了。 37 各人頭上光禿，鬍鬚剪短，手有劃傷，腰束麻布。 38 在摩押的各房頂上和街市上處處有人哀哭，因我打碎摩押，好像打碎無人喜悅的器皿。」這是耶和華說的。 39 「摩押何等毀壞，何等哀號，何等羞愧轉背！這樣，摩押必令四圍的人嗤笑、驚駭。」 40 耶和華如此說：「仇敵必如大鷹飛起，展開翅膀攻擊摩押。 41 加略被攻取，保障也被占據，到那日，摩押的勇士心中疼痛如臨產的婦人。 42 摩押必被毀滅，不再成國，因她向耶和華誇大。」 43 耶和華說：「摩押的居民哪，恐懼、陷坑、網羅都臨近你。 44 躲避恐懼的必墜入陷坑，從陷坑上來的必被網羅纏住，因我必使追討之年臨到摩押。」這是耶和華說的。

\subsubsection{後日摩押必復振興 48:45-47}
「躲避的人無力站在希實本的影下，因為有火從希實本發出，有火焰出於西宏的城，燒盡摩押的角和鬨嚷人的頭頂。 46 摩押啊，你有禍了！屬基抹的民滅亡了！因你的眾子都被擄去，你的眾女也被擄去。」 47 耶和華說：「到末後，我還要使被擄的摩押人歸回。」摩押受審判的話到此為止。

\subsection{以西結書-摩押亞捫亦必遭報 25:8-11}
%\subsubsection{摩押西珥亦必遭報 25:8-11}
主耶和華如此說：因摩押和西珥人說『看哪，猶大家與列國無異』， 9 所以我要破開摩押邊界上的城邑，就是摩押人看為本國之榮耀的伯耶西末、巴力免、基列亭， 10 好使東方人來攻擊亞捫人。我必將亞捫人之地交給他們為業，使亞捫人在列國中不再被記念。 11 我必向摩押施行審判，他們就知道我是耶和華。


\subsection{阿摩司書 2:1-3}
耶和華如此說：「摩押三番四次地犯罪，我必不免去他的刑罰，因為他將以東王的骸骨焚燒成灰。 2 我卻要降火在摩押，燒滅加略的宮殿，摩押必在鬨嚷呐喊吹角之中死亡。 3 我必剪除摩押中的審判者，將其中的一切首領和他一同殺戮。」這是耶和華說的。

\subsection{西番雅書 2:8-11}
我聽見摩押人的毀謗和亞捫人的辱罵，就是毀謗我的百姓，自誇自大，侵犯他們的境界。」 9 萬軍之耶和華以色列的神說：「我指著我的永生起誓：摩押必像所多瑪，亞捫人必像蛾摩拉，都變為刺草、鹽坑、永遠荒廢之地。我百姓所剩下的必擄掠他們，我國中所餘剩的必得著他們的地。」 10 這事臨到他們，是因他們驕傲，自誇自大，毀謗萬軍之耶和華的百姓。 11 耶和華必向他們顯可畏之威，因他必叫世上的諸神瘦弱，列國海島的居民各在自己的地方敬拜他。

